% SGA 3，第 VII_B 篇，英文重建稿——第 1 轮。
% 范围：权威底本分篇第 47 页，第 2.3 节至命题 2.3.1 之前的段落。
% 合并阅读本页码： 560.
% 再版印刷页码： 549.
% 本范围内原版页边定位： 516.
% 权威底本： Polo--Gille Exp7B-13oct24.pdf,
% SHA-256 D303690E0195654BBE7C3DAA5D89EA23E24CD709A3575217234AC19AE8836694.
% 曾参考 Jacob Reinhold 的 CC BY 4.0 Markdown 作为比对谱系。
% 本范围内没有图表；全部数学均为原生 LaTeX。
\subsection*{2.3}\phantomsection \label{sga3:VIIB:2.3}\addcontentsline{toc}{subsection}{2.3. 协变代数}

现在设 \(G\) 为任意形式 \(k\)-群，\footnote{\label{sga3:VIIB:nde:85}即 \(G\) 在 \(k\) 上未必拓扑平坦。}仿射代数为 \(A\)。仍以 \[
  \mathbf H(G)=\mathbf V_k^f(A)
\] 表示与 \(A\) 对偶的 \(\mathbf O_k\)-模，并以 \[
  \phi_G:
  \mathbf H(G)\otimes_k\mathbf H(G)
  \longrightarrow
  \mathbf H(G\times G)
\] 表示对 \(\operatorname{Alf}/k\) 的每个对象 \(C\) 由 \SGARefLink{sga3:VIIB:0.3.6}{0.3.6} 中的自然映射诱导的函子性同态：\[
\begin{aligned}
  &(A\widehat\otimes_kC)^\dagger
  \otimes_C
  (A\widehat\otimes_kC)^\dagger
  \\
  &\hspace{3em}\longrightarrow
  \left(
    (A\widehat\otimes_kC)
    \widehat\otimes_C
    (A\widehat\otimes_kC)
  \right)^\dagger.
\end{aligned}
\]。

若 \[
  m:G\times G\longrightarrow G
\] 为 \(G\) 的乘法，则复合 \[
  \mathbf H(G)\otimes_k\mathbf H(G)
  \xrightarrow{\ \phi_G\ }
  \mathbf H(G\times G)
  \xrightarrow{\ \mathbf H(m)\ }
  \mathbf H(G)
\] 赋予 \(\mathbf H(G)\) 在 \(\mathbf O_k\) 上的代数结构。对每个 \(C\in\operatorname{Alf}/k\)，\[
  \mathbf H(G)(C)
  =
  (A\widehat\otimes_kC)^\dagger
\] 的单位元就是 \(A\widehat\otimes_kC\) 的增广；参见 \SGARefLink{sga3:VIIA:2.1}{2.1}。\footnote{\label{sga3:VIIB:nde:86}编者调整了下列句子的顺序。}

若 \(G\) 在 \(k\) 上不拓扑平坦，态射 \(\phi_G\) 未必是同构。因此 \(G\) 的对角 \(x\mapsto(x,x)\) 所诱导的态射 \[
  \delta_G:
  \mathbf H(G)\longrightarrow\mathbf H(G\times G)
\] 未必通过 \[
  \mathbf H(G)\otimes_k\mathbf H(G).
\] 分解。所以 \(\mathbf H(G)\) 未必是 \(\mathbf O_k\) 上的双代数。因此在一般情形中，仅称 \(\mathbf H(G)\) 为形式群 \(G\) 的\emph{协变代数}。

当然，当 \(G\) 在 \(k\) 上拓扑平坦时，\(\phi_G\) 为同构，因而重新得到 \SGARefLink{sga3:VIIB:2.2.1}{2.2.1} 中定义的 \(\mathbf H(G)\) 上的 \(\mathbf O_k\)-双代数结构。
